%! Inverses of Exponential Functions — Unit 5, Lesson 5.2 — Homework (Answer Key)
\documentclass[10pt]{article}
\usepackage{precalculus-article}
\usepackage{precalculus-key}   % pulls in -boxes; do NOT also load -boxes
\usepackage{grading}
\gradingassignment{precalc-u05-l02-hw}
\gradingcourse{precalculus}
\newcommand{\TallMath}[1]{$\displaystyle #1\rule[-1.4em]{0pt}{3.2em}$}

\begin{document}

\pageheader{Unit 5, Lesson 5.2}{Homework --- Answer Key}
\namedateperiod

\vspace{0.1in}

\begin{enumerate}[leftmargin=1.8em, itemsep=14pt]

\item Let $g(x) = 6^x$.

\begin{enumerate}[leftmargin=1.5em, itemsep=6pt, label=\alph*.]
  \item Write the inverse function $f(x)$.

  \vspace{4pt}
  $f(x) = $ \smash{\gradedblank[id=hw1a, type=short, points=1, answer={\log_6 x}]{4cm}{13pt}}

  \item Complete the table and inverse pairs.

  \vspace{4pt}
  \begin{center}
  \renewcommand{\arraystretch}{1.8}
  \begin{tabular}{|c|c|c|}
  \hline
  \rowcolor{sky}
  $x$ & $g(x) = 6^x$ & Point on $f$: $(g(x),\; x)$ \\
  \hline
  $0$ & \smash{\gradedblank[id=hw1b1, type=short, points=1, answer={1}]{1.5cm}{14pt}} & \smash{\gradedblank[id=hw1b2, type=short, check=text, points=1, answer={(1, 0)}, display={$(1,\;0)$}]{3cm}{14pt}} \\
  \hline
  $1$ & \smash{\gradedblank[id=hw1b3, type=short, points=1, answer={6}]{1.5cm}{14pt}} & \smash{\gradedblank[id=hw1b4, type=short, check=text, points=1, answer={(6, 1)}, display={$(6,\;1)$}]{3cm}{14pt}} \\
  \hline
  $2$ & \smash{\gradedblank[id=hw1b5, type=short, points=1, answer={36}]{1.5cm}{14pt}} & \smash{\gradedblank[id=hw1b6, type=short, check=text, points=1, answer={(36, 2)}, display={$(36,\;2)$}]{3cm}{14pt}} \\
  \hline
  \end{tabular}
  \end{center}
\end{enumerate}

\item Verify $f(x) = \log_5 x$ and $g(x) = 5^x$ are inverses.

\vspace{4pt}
$f(g(x)) = \log_5(5^x) = $ \smash{\gradedblank[id=hw2a, type=short, points=1, answer={x}]{3cm}{13pt}}

\vspace{6pt}
$g(f(x)) = 5^{\log_5 x} = $ \smash{\gradedblank[id=hw2b, type=short, points=1, answer={x}]{3cm}{13pt}}

\vspace{4pt}
Conclusion: \writebox{\smash{\gradedblank[id=hw2c, type=rubric, points=1,
  display={Both compositions equal $x$, so $f$ and $g$ are inverses.}]{431pt}{13pt}}}

\item Algae model $A(t) = 10^t$.

\begin{enumerate}[leftmargin=1.5em, itemsep=6pt, label=\alph*.]
  \item $t(A) = $ \smash{\gradedblank[id=hw3a, type=short, points=1, answer={\log_{10} A}]{4cm}{13pt}}

  \item $t(1000) = \log_{10} 1000 = $ \smash{\gradedblank[id=hw3b1, type=short, points=1, answer={3}]{3cm}{13pt}}

  \vspace{4pt}
  In context: \writebox{\smash{\gradedblank[id=hw3b2, type=rubric, points=1,
  display={After 3 days the algae covers 1000 cm$^2$.}]{419pt}{13pt}}}
\end{enumerate}

\item \textbf{Conceptual.}

\writegap\\
\writegap\\
\writebox{\gradedblank[id=hw4, type=rubric, points=2,
  display={Exponential: adding to the input multiplies the output by $b$. Logarithm: multiplying the input by $b$ adds 1 to the output. Inverses swap inputs and outputs, so the patterns are reversed.}]{\linewidth}{41pt}}\\

\item Let $h(x) = 2\log_3 x$.

\begin{enumerate}[leftmargin=1.5em, itemsep=6pt, label=\alph*.]
  \item $b = $ \smash{\gradedblank[id=hw5a1, type=short, points=1, answer={3}]{2cm}{13pt}} \qquad $a = $ \smash{\gradedblank[id=hw5a2, type=short, points=1, answer={2}]{2cm}{13pt}}

  \item $h(9) = 2\log_3 9 = 2 \cdot 2 = $ \smash{\gradedblank[id=hw5b1, type=short, points=1, answer={4}]{3cm}{13pt}} \qquad
        $h(27) = 2\log_3 27 = 2 \cdot 3 = $ \smash{\gradedblank[id=hw5b2, type=short, points=1, answer={6}]{3cm}{13pt}}

  \item Compare the outputs of $h(x)$ and $\log_3 x$; describe the dilation.

  \writebox{\smash{\gradedblank[id=hw5c, type=rubric, points=2,
  display={Each output of $h$ is twice that of $\log_3 x$: a vertical stretch by 2.}]{\linewidth}{13pt}}}
\end{enumerate}

\item \textbf{AP-style MC.}

\begin{enumerate}[leftmargin=2em, itemsep=4pt, label=(\Alph*)]
  \item $g(x) = x^4$
  \item $g(x) = 4x$
  \item $g(x) = 4^x$ \quad \textcolor{keyred}{\textbf{$\leftarrow$ correct}}
  \item $g(x) = \log x$ \hfill\textit{Letter:}~\smash{\gradedblank[id=hw6, type=choice, points=1, answer=C]{0.42in}{13pt}}
\end{enumerate}

\begin{teachernote}
EK 2.10.A.4: $f(x) = \log_b x$ is the reflection of $g(x) = b^x$ over $y = x$, so the
inverse of $\log_4 x$ is $4^x$. (A) is a power function, not an exponential; (B) is linear;
(D) has a different base (base 10).
\end{teachernote}

\end{enumerate}

\vspace{0.12in}

\begin{reflectionbox}
What is one thing from today's lesson that you feel confident about, and one question you
still have about inverse functions or logarithms?

\writelines{2}
\end{reflectionbox}

\end{document}
